\chapter*{Conclusions}
\addcontentsline{toc}{chapter}{Conclusions}
\label{ch:conclusions}
\markboth{Conclusions}{Conclusions}

\section{Summary of Approach}

\indent This dissertation investigated information security in electronic banking and electronic finance service systems, with emphasis on elliptic curve cryptography for authentication, key exchange, and data protection. A systematic threat analysis (Chapter 1) established the attack landscape and identified critical limitations in deployed authentication mechanisms (Section 1.2) and data-security mechanisms (Section 1.3) \hyperref[ref:11]{[11]}\hyperref[ref:34]{[34]}\hyperref[ref:35]{[35]}.

\indent Two novel ECC-based methodologies were proposed, analyzed, implemented, and evaluated: ECC-DTB-AKA for client-server authenticated key agreement with device and transaction binding (Chapter 2), and ECC-HISE for server-server hierarchical secure envelopes with Merkle audit chains (Chapter 3). Chapter 4 demonstrated integrated architecture and programmatic implementation. Chapter 5 evaluated effectiveness against network attacks, reviewed prior work, and synthesized vulnerability findings \hyperref[ref:7]{[7]}\hyperref[ref:8]{[8]}\hyperref[ref:30]{[30]}.

\section{Contributions}

\indent The principal contributions of this research are:

\indent First, the design and formal analysis of ECC-DTB-AKA, a three-round protocol achieving mutual authentication, forward secrecy, device binding, and per-transaction key derivation using SECP256R1 ECDH, ECDSA, and HKDF \hyperref[ref:5]{[5]}\hyperref[ref:6]{[6]}\hyperref[ref:9]{[9]}.

\indent Second, the design and formal analysis of ECC-HISE, a hierarchical envelope scheme combining domain-separated key derivation, AES-256-GCM encryption, Ed25519 signatures, and Merkle-tree batch integrity for inter-bank settlement \hyperref[ref:2]{[2]}\hyperref[ref:10]{[10]}\hyperref[ref:30]{[30]}.

\indent Third, a dual-layer integrated security architecture unifying both schemes in an e-financial service system with defined trust boundaries and component interactions.

\indent Fourth, a working software prototype with REST APIs, benchmark suites, and end-to-end integration tests providing reproducible quantitative evaluation.

\indent Fifth, comprehensive prior-work analysis positioning the proposed schemes relative to TLS 1.3, OAuth 2.0 PKCE, FIDO2, ECIES, static PSK, and blockchain audit approaches \hyperref[ref:1]{[1]}\hyperref[ref:28]{[28]}\hyperref[ref:13]{[13]}.

\section{Limitations}

\indent This research has limitations that should be acknowledged. The prototype uses software key storage rather than hardware security modules. Formal security proofs are presented as structured proof sketches rather than machine-verified derivations. Certificate management is simplified compared to production PKI. The settlement demonstration performs sender and receiver verification in a single process. These limitations do not invalidate the cryptographic designs but indicate areas requiring hardening for production deployment \hyperref[ref:4]{[4]}.

\section{Future Work}

\indent Future research directions include: integration of NIST post-quantum algorithms (ML-KEM, ML-DSA) for long-term confidentiality; HSM-backed key storage and FIPS 140-3 validation; formal verification using ProVerif or Tamarin provers; field trials with partner financial institutions; meta-Merkle hierarchies for million-record national settlement; and composition with FIDO2/WebAuthn for phishing-resistant client key enrollment \hyperref[ref:15]{[15]}\hyperref[ref:30]{[30]}\hyperref[ref:54]{[54]}.

